\subsection{A 引用格式}\label{a-citation-reference}

\begin{Shaded}
\begin{Highlighting}[]
\CommentTok{引用本工作时，请使用以下 BIBTEX 条目：}
\VariableTok{@techreport}\NormalTok{\{}\OtherTok{mai\_thinking\_1}\NormalTok{,}
\DataTypeTok{title}\NormalTok{=\{MAI{-}Thinking{-}1: Building a Hill{-}Climbing Machine\},}
\DataTypeTok{author}\NormalTok{=\{The Microsoft AI Team\},}
\DataTypeTok{institution}\NormalTok{=\{Microsoft AI\},}
\DataTypeTok{year}\NormalTok{=\{2026\},}
\DataTypeTok{urldate}\NormalTok{=\{2026{-}06{-}02\},}
\DataTypeTok{url}\NormalTok{=\{https://microsoft.ai/pdf/mai{-}thinking{-}1.pdf\}}
\NormalTok{\}}
\end{Highlighting}
\end{Shaded}

\subsection{B 预训练数据管线细节}\label{b-pretraining-data-pipeline-details}

本附录进一步介绍第 2.4 节所述主要预训练数据源的处理细节。

\subsection{B.1 网页 HTML}\label{b.1-web-html}

我们的网页 HTML 语料大部分来自内部抓取。在完成初始的页面发现与筛选后，我们抓取并解析了约 1.2 万亿个页面。HTML 使用 Trafilatura（Barbaresi, 2021）转换为文本，输出抽取出的文本以及语言标识和相应的置信度分数。在抽取过程中，我们推断文档编码，而不是在处理前假定其为 UTF-8。这种方式在处理带有旧式、缺失或错误声明编码的多语言网页数据时更具鲁棒性。

解析完成后，所有页面都要经过政策合规、成人内容、与网页 PDF 语料重叠以及域名黑名单等过滤。除第 2.4 节所述的微软标准政策外，我们还应用 UT1 黑名单（Prigent, 2026）来移除成人内容与盗版相关域名。总体而言，这些过滤将语料从 1.2 万亿页缩减至 7940 亿页。鉴于网络上 AI 生成内容日益普遍，我们还使用内部 AI 内容检测模型对页面评分，并通过人工检查识别包含大量 AI 生成内容的域名；这些域名将从训练语料中过滤掉。

过滤后的网页语料随后进行精确去重与模糊去重。我们先通过转小写和去除首尾空白来规范化文档，然后使用 MD5 哈希移除完全重复项。这一步将语料从 7940 亿缩减至 4230 亿篇文档。我们的模糊去重采用第 2.4.3 节所述的 LSH 方法与参数。模糊去重后，内部抓取语料包含 734 亿篇英文文档和 1165 亿篇非英文文档。

我们对 Common Crawl 使用相同的管线。从约 3000 亿个页面出发，我们为每个 URL 保留最新抓取的内容，将语料缩减至约 1000 亿个页面。经过过滤、去重、与内部网页语料合并，以及最后一轮基于精确 URL 和内容级的模糊去重后，Common Crawl 部分包含 242 亿个页面。

在下游处理中，我们使用 Qwen3-Embedding-0.6B（Yang et al., 2025）为所有网页计算文本嵌入。得到的语料随后通过面向通用网页、STEM 页面、代码页面和关键领域的专属管线进行处理。

通用网页。对于通用网页，我们结合过滤与质量分箱，以便对这一大规模语料进行更细粒度的数据配比。首先，我们将属性模型应用于每篇文档的文本嵌入。这些属性模型对与预训练相关的属性进行评分，如教育价值、事实准确性与可靠性、写作质量、信息密度、推理内容以及通识知识价值。在这些属性之上，我们训练一个质量模型，并过滤掉底部 70\% 的英文文档。随后应用 Gopher 风格（Rae et al., 2022）的启发式过滤器，其基于词统计、标点模式、常见 n-gram 及相关文档质量特征。这一步从内部抓取语料中保留约 46 亿篇英文文档，从英文 Common Crawl 中保留 28 亿篇。为便于数据配比，这些文档结合元数据与基于嵌入的质量评分进行分箱。多语言网页数据的处理方式与英文通用网页数据类似，但采用语言感知的过滤与质量评分。

除广泛的通用网页管线外，我们还构建了一个高质量的人工策展网页内容子集，以更好地捕捉引人入胜、文笔上乘的长文内容。我们从人工策展渠道（如杂志和长文阅读聚合站）以及社交平台上的外部引用链接获取候选域名，并从网页语料中取回相应页面。

STEM 页面。预训练阶段对高质量 STEM 内容的广泛覆盖，对于下游的数学与科学推理能力至关重要。然而，这类内容在网上相对稀少，且通常包含公式、代码片段和技术标记等结构化元素，通用网页处理管线难以很好地处理这些内容。因此，我们构建了一条专用管线来识别并抽取网页中的 STEM 内容。

我们使用基于嵌入的分类器从三个维度识别候选 STEM 页面：主题、教育价值和教育水平。主题由七个二分类器建模，分别对应数学、物理、统计、化学、生物、工程和计算机科学。一个共享的教育价值分类器用于将教育类内容（如研究论文、Q\&A 论坛、教程、讲义和博客）与低价值页面（如产品列表、引文索引、商业内容和模板化页面）区分开。最后，一个教育水平分类器将每篇文档划分为高中、本科或研究生水平，用于数据配比。为提高语料质量，我们还应用了一套针对 STEM 内容定制的启发式规则，因为通用网页过滤器往往会过度惩罚包含公式和数值表达式的文档。这些启发式规则会移除 PDF URL、计算器域名、专利域名、已知 AI 生成内容域名以及被判定为低质量的域名，并根据词数、数字占比、字母占比、重复行比例、停用词密度以及以省略号结尾的行占比来过滤文档。

由于 Trafilatura 并不总能很好地保留数学公式和代码，我们为 STEM 文档构建了定制的抽取管线。从原始 HTML 出发，我们对 MathML 和 LATEX 内容进行规范化，并将文档转换为 Markdown。随后将文档切分为多个小节，每个公式和代码块作为独立单元保留，并使用语言模型将每个小节分类为保留或移除。该模型仅限于二元取舍决策，不能生成合成内容。最后，我们使用 OpenWebMath 风格的后处理（Gokaslan and Cohen, 2019）对抽取出的内容进行清洗和规范化，使公式的定界保持一致。

该管线产出一个 680B tokens 的英文 STEM 语料，其中 76B 与数学相关。对于多语言数据，我们额外获得 760B 的 STEM tokens，其中包含 58B 高质量且与数学相关的 tokens。

代码页面。对于代码类网页，我们复用 STEM 页面管线，包括主题分类、教育价值分类和基于 LLM 的抽取，并聚焦于计算机科学主题进行过滤。为进一步提升质量，我们使用 Qwen3-30B（Yang et al., 2025）对每篇候选文档进行评分。评判提示（judge prompt）使用 GEPA / DSPy（Agrawal et al., 2026）配合约 2,000 个人工标注进行优化。在通过额外启发式规则过滤掉低质量文档后，我们获得约 233B tokens 的数据集。

关键领域。对于选定的高价值领域，我们使用改进的抽取管线，以更好地保留领域特有的结构与内容。例如，对于 Wikipedia，我们在收录上宁偏宽松，避免在标准抽取可能丢弃重要页面元素时遗漏有用信息。

\subsection{B.2 网页 PDF}\label{b.2-web-pdfs}

我们收集了一个网络抓取的 PDF 语料，约含 10B 篇文档，涵盖教育、生活、技术和专业文档。基于从 PDF 字节中抽取出的文本，我们使用启发式规则和分类器将其过滤至 620M 篇文档，供进一步处理。为将文档转换为可训练文本，我们使用 Azure Document Intelligence（Azure DI）。

我们对 OCR 输出应用若干后处理步骤：

• 公式与表格规范化

• 移除模板化内容，如图眉、页脚和图中检测到的文本

• 移除 OCR 伪影，如字形错误

• 行展开：去除换行符和连字符

• 移除学术文档中的参考文献部分

我们使用文本级分类器和 PDF 元数据启发式规则来移除 SEO 和垃圾内容。特别是，PDF 元数据字段 ``creator'' 和 ``producer'' 是过滤时高精度的有效信号。我们进一步使用文本长度、gzip 可压缩性、字母数字占比和表格占比等启发式规则以及质量分类器对文档进行过滤。

经过模糊去重后，PDF 数据集包含 1.8T 英文 tokens 和 1.85T 多语言 tokens。为支持数据配比，我们将文档分类为教育类或非教育类。教育类文档进一步划分为数学、计算机科学、STEM 和非 STEM，并指定为高中、本科和研究生中的某一教育水平。

\subsection{B.3 图书与期刊}\label{b.3-books-and-journals}

我们从多家供应商处获取图书和期刊，包括通过与出版商的直接协议。针对每家供应商，我们构建一条专用的摄取管线，以适配其底层数据格式，并与内容附带的用途限制保持一致。每条管线将原始内容转换为规范的文本格式。标准化后的文本随后经过 OCR 伪影清洗、启发式过滤和内容筛选。清洗后的文本经过精确去重、标题去重和模糊去重，以减少不必要的重复。最后，使用 LLM 进行标注，为语料中的每篇文档赋予主题和质量标签，以评估每部作品对 LLM 预训练的用处，并更好地控制数据配比。

\subsection{B.4 公开 GitHub}\label{b.4-public-github}

我们 7.4T-token 的代码预训练语料大部分来自公开 GitHub 仓库，并组织为三个数据集：files、commits 和 pull requests（PRs）。三个数据集均经过共享的过滤、去重、去污染和质量评分管线。

共享处理。我们首先应用启发式文件级过滤器，以移除非代码和低质量数据。这些过滤器包括：

• 移除 node\_modules、build、\_\_pycache\_\_ 和 .vscode 等垃圾文件夹中的文件。

• 通过路径模式检测生成的代码，如 \_pb2.py 之类的 protobuf 输出以及 .d.ts 之类的 TypeScript 声明文件。

• 排除非代码文件类型，如二进制文件、图像、压缩包、文档和字体。

• 拒绝超过 30K 字符或包含过长行的文件。

• 字符构成检查，用于过滤掉数据转储。

随后我们通过 SHA-512 哈希进行精确去重，再进行基于 MinHashLSH 的模糊去重，以及使用 Qwen3-Embedding-0.6B（Yang et al., 2025）的 1024 维嵌入上的余弦相似度进行语义去重。我们还会针对推理训练和评测中使用的编程题进行去污染。最后，每个样本被赋予一个质量分数，并放入离散的质量分箱，以支持质量感知采样。

在完成各数据集特有的处理后，每个数据集中的样本被贪婪地打包为定长序列，files、commits 和 PRs 分别产生约 1.26T、4.5T 和 1.19T tokens。

Files。原始 files 数据集包含所有仓库中每个文件的最新版本。我们使用与 StarCoder2（Lozhkov et al., 2024）中所述类似的过滤器进行额外的启发式过滤，并特别丢弃 token 数量较大的仓库中的非编程文件。文件按仓库分组，按仓库结构的深度优先遍历顺序排序，并拼接为仓库级序列。

Commits。原始 commits 数据集包含从每个仓库抽取的最新的 10K 条提交。每一行包含所有被修改文件的提交前状态以及该提交所应用的补丁。作为某条 PR 组成部分的提交会被移除，以避免跨数据集冗余。被修改文件的提交前状态用作 loss-masked 前缀，而可训练文本仅由提交补丁组成，以 git-diff 或 search-and-replace 格式呈现。

Pull requests。原始 PR 数据集通过从克隆的仓库中抽取所有组成提交并为每条 PR 关联相关元数据而构建。每一行包含 PR 标题、正文、关联的 issue、issue 评论、reviews、threads、review 评论、提交前文件状态以及按顺序排列的组成提交列表。PR 数据集还额外针对 SWE-bench Verified 进行去污染，移除基准测试中使用的所有 PR。被修改文件的提交前状态用作 loss-masked 前缀；可训练文本包含其余字段，包括 PR 标题、正文、关联的 issue、reviews、评论以及各组成提交的补丁。当元数据可映射到特定提交（如 review 评论和 threads）时，它会被紧跟在相应提交补丁之后放置，而非放在 PR 级别。我们还对过长的前缀应用前缀压缩，随机选择每条 git 补丁位置周围的行，以适应固定的 token 预算。压缩后仍超过 256K tokens 的序列将被丢弃。
